\scanpage{76}
\mspage{- 3 -}
% A long pen stroke crosses b), c), section 3 and the footnotes.
\deleted{b)\quad Les travaux de Igusa sur les ``vanishing cycles'' (dont il ne sera question que dans un séminaire ultérieur), et la démonstration par Igusa de \uline{l'inégalité de Picard} $\rho\handwritten{\leq}b_2$, en toute caractéristique, que nous obtenons ici de façon assez formelle (en même temps que les théorèmes de finitude de Néron-Sévéri), une fois en possession d'une théorie cohomologique explicite.}

\deleted{c)\quad Les résultats de \uline{Ogg} et \uline{Chafarevitch} sur la cohomologie des corps de fonction d'une variable à coefficient \deleted{des}\handwritten{\inserted{des}} variétés abéliennes. Ces auteurs obtiennent en particulier un équivalent du \uline{théorème de dualité} pour les courbes algébriques\handwritten{,} en même temps qu'une formule des caractéristiques d'Euler-Poincaré. Pour un exposé en termes de cohomologie étale, cf \handwritten{[ \phantom{0} ]}.}

\deleted{3.\qquad Les exposés I \unsure{à|-- à} VI. reprennent la théorie des faisceaux du point de vue des sites, ou mieux, des topos, en complétant le séminaire Artin 1962 sur de nombreux points. La présentation\handwritten{,} et les résultats non contenus dans Artin\handwritten{,} sont dûs pour l'essentiel à J. Giraud et J.L. Verdier.}
\margin{\handwritten{$\times$}}

\deleted{Les exposés VII \handwritten{et VIII} développent les premières définitions et propriétés relatives à la topologie étale et la cohomologie étale.}

\deleted{Les exposés \handwritten{IX}, \handwritten{X} donnent les premiers résultats quantitatifs élémentaires relatifs aux coefficients de torsion\handwritten{,} concernant en particulier la cohomologie des courbes algébriques, la dimension cohomologique, la notion de faisceaux constructibles.}

% slanted
\margin{\handwritten{X.3.10\quad \uline{XII}}}
\deleted{Les résultats des exposés XII à XVI\textsuperscript{(2)} sont centrés autour des deux \handwritten{``}théorèmes de changement de base'', \handwritten{\deleted{Nous avons pris soin de donner}\inserted{qui techniquement constituent le résultat central de ce séminaire.}}}

% slanted
\margin{\handwritten{\unsure{X.3.4|X.34}}}
\deleted{Les exposés XVII à \handwritten{\deleted{\unsure{XIX|XVIII}}\unsure{XVIII|XIX}}\textsuperscript{(3)} présentent la théorie de la cohomologie ``à supports compacts'' et le \uline{théorème de\deleted{s} dualité\deleted{s} globale\deleted{s}}, qui impliquent la \uline{dualité de Poincaré} pour les variétés algébriques non singulières.}

\deleted{\textsuperscript{(2)} Obtenus entre Septembre 1962 et Mars 1963 par M. Artin et A. Grothendieck}

\deleted{\textsuperscript{(3)} Obtenus en Février-Mars 1963 par A. Grothendieck, en même temps que le ``théorème de dualité locale'' qui sera exposé \handwritten{dans \deleted{\unsure{1965|SGA}} SGA 1965}}
% archive sheet number: 75
